\documentclass[10pt]{article}

% COMP 2300 Proofs Core companion: student-facing reference and practice.
% The longer proofs_cheat_sheet.tex remains the extended reference.
\usepackage[letterpaper,margin=0.56in,includehead,includefoot,
  headheight=15pt,headsep=7pt,footskip=20pt]{geometry}
\usepackage[T1]{fontenc}
\usepackage{amsmath,amssymb,mathtools}
\usepackage{microtype}
\usepackage[dvipsnames,svgnames]{xcolor}
\usepackage{multicol}
\usepackage{enumitem}
\usepackage{tabularx,booktabs,array}
\usepackage[most]{tcolorbox}
\usepackage{fancyhdr}
\usepackage{titlesec}
\usepackage[hidelinks]{hyperref}

\definecolor{CoreNavy}{HTML}{18364A}
\definecolor{CoreTeal}{HTML}{176B6B}
\definecolor{CoreGold}{HTML}{9A6500}
\definecolor{CoreRed}{HTML}{963D38}
\definecolor{CoreBlueTint}{HTML}{EDF4F8}
\definecolor{CoreTealTint}{HTML}{EDF8F6}
\definecolor{CoreGoldTint}{HTML}{FFF7E5}
\definecolor{CoreRedTint}{HTML}{FCEFED}

\setlength{\parindent}{0pt}
\setlength{\parskip}{2.5pt}
\setlength{\columnsep}{0.25in}
\setlength{\columnseprule}{0.25pt}
\setlength{\emergencystretch}{1.2em}
\setlength{\abovedisplayskip}{3pt}
\setlength{\belowdisplayskip}{3pt}
\setlist{leftmargin=*,itemsep=1pt,topsep=1.5pt,parsep=0pt,partopsep=0pt}
\raggedbottom

\titleformat{\section}{\Large\bfseries\color{CoreNavy}}{}{0pt}{}
  [\vspace{-3pt}\color{CoreNavy}\rule{\linewidth}{0.7pt}]
\titlespacing*{\section}{0pt}{3pt}{4pt}

\pagestyle{fancy}
\fancyhf{}
\fancyhead[L]{\textbf{COMP 2300}}
\fancyhead[C]{\textbf{Proofs Core Companion}}
\fancyhead[R]{Definitions and Worked Examples}
\fancyfoot[L]{Core sections A--E}
\fancyfoot[C]{\thepage}
\fancyfoot[R]{Study reference---not used during a Core}
\renewcommand{\headrulewidth}{0.5pt}
\renewcommand{\footrulewidth}{0.3pt}

\tcbset{enhanced,before skip=2.5pt,after skip=2.5pt,boxrule=0.55pt,arc=1.2mm,
  left=4pt,right=4pt,top=3pt,bottom=3pt,fonttitle=\bfseries}
\newtcolorbox{definitionbox}[1]{title={#1},colback=CoreBlueTint,
  colframe=CoreNavy,colbacktitle=CoreNavy,coltitle=white}
\newtcolorbox{methodbox}[1]{title={#1},colback=CoreTealTint,
  colframe=CoreTeal,colbacktitle=CoreTeal,coltitle=white}
\newtcolorbox{workedbox}[1]{title={#1},colback=white,
  colframe=CoreNavy,colbacktitle=CoreBlueTint,coltitle=CoreNavy}
\newtcolorbox{warningbox}[1]{title={#1},colback=CoreRedTint,
  colframe=CoreRed,colbacktitle=CoreRed,coltitle=white}
\newtcolorbox{checklistbox}[1]{title={#1},colback=CoreGoldTint,
  colframe=CoreGold,colbacktitle=CoreGold,coltitle=white}
\newcommand{\term}[1]{\textbf{\color{CoreNavy}#1}}

\begin{document}
\small

\section*{A. Inference and Proof Structure}
\begin{multicols}{2}
\begin{definitionbox}{Argument, validity, and countermodel}
An \term{argument} consists of premises and a conclusion. It is \term{valid}
when every valuation or interpretation that makes all premises true also makes
the conclusion true. A \term{countermodel} makes every premise true and the
conclusion false; one countermodel proves invalidity.
\end{definitionbox}

\begin{definitionbox}{Propositional rules of inference}
\begin{align*}
p,\ p\to q&\therefore q &&\text{Modus Ponens (MP)}\\
\neg q,\ p\to q&\therefore\neg p &&\text{Modus Tollens (MT)}\\
p\to q,\ q\to r&\therefore p\to r &&\text{Hypothetical Syllogism}\\
p\lor q,\ \neg p&\therefore q &&\text{Disjunctive Syllogism}\\
p,\ q&\therefore p\land q &&\text{Conjunction}\\
p\land q&\therefore p &&\text{Simplification}\\
p&\therefore p\lor q &&\text{Addition}\\
p\lor q,\ p\to r,\ q\to r&\therefore r &&\text{Proof by cases}
\end{align*}
\end{definitionbox}

\begin{workedbox}{A propositional derivation}
\textbf{Target:} derive $\neg p$ from $p\to q$, $q\to r$, and $\neg r$.
Start at $\neg r$ and move backward through the conditional chain:
\[
\begin{array}{rll}
1.&p\to q&\text{premise}\\
2.&q\to r&\text{premise}\\
3.&\neg r&\text{premise}\\
4.&\neg q&\text{MT, 2 and 3}\\
5.&\neg p&\text{MT, 1 and 4.}
\end{array}
\]
Line 4 removes the possibility of $q$; line 5 then removes the possibility of
$p$. Name both the rule and the lines used.
\end{workedbox}

\columnbreak

\begin{definitionbox}{Quantifier rules}
\term{Universal instantiation (UI):} from $\forall x\,P(x)$ infer $P(c)$ for
any named object $c$ in the domain.

\term{Universal generalization (UG):} after proving $P(c)$ for a completely
arbitrary $c$, infer $\forall x\,P(x)$.

\term{Existential instantiation (EI):} from $\exists x\,P(x)$ introduce a
fresh name $a$ with $P(a)$. It denotes some witness, not an arbitrary object.

\term{Existential generalization (EG):} from $P(c)$ infer $\exists x\,P(x)$.
\end{definitionbox}

\begin{methodbox}{Universal and existential witnesses differ}
\textbf{Universal route:} from $\forall x(P(x)\to Q(x))$ and $\forall xP(x)$,
(1) let $c$ be arbitrary; (2) UI gives $P(c)\to Q(c)$ and $P(c)$; (3) MP
gives $Q(c)$; (4) because $c$ was arbitrary, UG gives $\forall xQ(x)$.

\textbf{Existential warning:} from $\exists xP(x)$ and $\exists xQ(x)$, EI may
introduce different witnesses $a$ and $b$. You may not conclude
$P(a)\land Q(a)$ or
$\exists x(P(x)\land Q(x))$ without evidence that the witnesses coincide.
\end{methodbox}

\begin{workedbox}{A countermodel for reused witnesses}
\textbf{Build:} use domain $\{1,2\}$; let $P$ hold only of $1$ and $Q$ only of
$2$. \textbf{Check premises:} $1$ witnesses $\exists xP(x)$, and $2$ witnesses
$\exists xQ(x)$. \textbf{Check conclusion:} neither object satisfies both
predicates, so $\exists x(P(x)\land Q(x))$ is false. True premises and a false
conclusion prove invalidity.
\end{workedbox}

\begin{warningbox}{Invalid forms}
$p\to q,q\therefore p$ affirms the consequent. $p\to q,\neg p\therefore
\neg q$ denies the antecedent. Neither is valid. A conclusion may be true even
when the displayed argument for it is invalid.
\end{warningbox}

\begin{checklistbox}{A readiness}
Can you execute a short derivation with named rules, recognize invalid forms,
use UI/EI with correct scopes, and construct a small countermodel when an
existential witness has been reused?
\end{checklistbox}
\end{multicols}

\clearpage
\section*{B. Direct Proofs from Definitions}
\begin{multicols}{2}
\begin{definitionbox}{Definitions used on this Core}
For integers $n$:
\[
\begin{array}{ll}
n\text{ is even}&\Longleftrightarrow n=2k\text{ for some }k\in\mathbb Z,\\
n\text{ is odd}&\Longleftrightarrow n=2k+1\text{ for some }k\in\mathbb Z,\\
d\mid n&\Longleftrightarrow n=dk\text{ for some }k\in\mathbb Z.
\end{array}
\]
A real $r$ is \term{rational} iff $r=a/b$ for integers $a,b$ with $b\ne0$;
otherwise it is \term{irrational}. $A\subseteq B$ means
$\forall x(x\in A\to x\in B)$.
\end{definitionbox}

\begin{methodbox}{Direct-proof template}
For ``if $P$, then $Q$'':
\begin{enumerate}
\item Let the objects be arbitrary subject to the stated domain.
\item Assume $P$.
\item Unpack $P$ using definitions and introduce correctly scoped witnesses.
\item Calculate or reason to the defining form of $Q$.
\item State why the new witness belongs to the required domain.
\end{enumerate}
\end{methodbox}

\begin{workedbox}{B: an even-square proof}
\textbf{Set up:} let arbitrary $n\in\mathbb Z$ be even, so $n=2k$ for some
$k\in\mathbb Z$. \textbf{Calculate and expose a factor of $2$:}
\[
n^2=(2k)^2=4k^2=2(2k^2).
\]
\textbf{Close:} $2k^2\in\mathbb Z$, so the display has form
$2(\text{integer})$. Thus $n^2$ is even. The new witness is $2k^2$.
\end{workedbox}

\begin{workedbox}{B: two odd integers}
\textbf{Set up:} let arbitrary odd integers $a,b$ satisfy $a=2j+1$ and
$b=2k+1$ for some $j,k\in\mathbb Z$. Use different witnesses. \textbf{Combine:}
\[
a+b=2j+1+2k+1=2(j+k+1).
\]
\textbf{Close:} $j+k+1\in\mathbb Z$, so $a+b$ is twice an integer and is even.
\end{workedbox}

\columnbreak

\begin{workedbox}{B: sum of rational numbers}
\textbf{Set up:} let $r,s\in\mathbb Q$. Choose integers $a,b,c,d$ with
$b,d\ne0$, $r=a/b$, and $s=c/d$. \textbf{Use a common denominator:}
\[
r+s=\frac{ad+bc}{bd}.
\]
\textbf{Close:} $ad+bc,bd\in\mathbb Z$, and $bd\ne0$. The displayed fraction
satisfies the definition of rational; lowest terms are not required.
\end{workedbox}

\begin{workedbox}{B: direct element proof}
Assume $A\subseteq B$. \textbf{Start on the left:} let arbitrary
$x\in A\cap C$. \textbf{Unpack:} $x\in A$ and $x\in C$. \textbf{Use the
hypothesis:} $x\in A$ and $A\subseteq B$ give $x\in B$. \textbf{Repack:}
$x\in B$ and $x\in C$ give $x\in B\cap C$. Since $x$ was arbitrary, the
subset claim follows.
\end{workedbox}

\begin{definitionbox}{Necessary proof language}
\term{Arbitrary} means no special property beyond the hypotheses was chosen.
A \term{witness} is a specific object whose existence establishes an
existential claim. A \term{lemma} is a proved supporting result. A
\term{counterexample} satisfies a universal claim's hypotheses but violates
its conclusion.
\end{definitionbox}

\begin{warningbox}{Common errors}
Do not write $a=2k$ and $b=2k$ merely because both are even; their witnesses
may differ. Do not divide by a variable without knowing it is nonzero. Do not
give an example when the claim is universal. End by connecting the algebra to
the exact definition requested.
\end{warningbox}

\begin{checklistbox}{B readiness}
Can you identify the target definition, introduce independent witnesses,
perform the central algebraic or element-chasing step, verify the new witness,
and state the conclusion for arbitrary inputs?
\end{checklistbox}
\end{multicols}

\clearpage
\section*{C--D. Indirect Methods, Cases, Existence, and Refutation}
\begin{multicols}{2}
\begin{definitionbox}{Contrapositive and contradiction}
The contrapositive of $P\to Q$ is $\neg Q\to\neg P$ and is logically
equivalent to the original. A \term{proof by contradiction} assumes the
negation of the desired conclusion and derives an impossibility such as
$R\land\neg R$. State which method you are using.
\end{definitionbox}

\begin{methodbox}{C: contraposition workflow}
Write the contrapositive explicitly, assume its antecedent, and prove its
consequent. If the antecedent is a disjunction, use cases. Do not assume the
original conclusion.
\end{methodbox}

\begin{workedbox}{C: product odd implies both factors odd}
\textbf{Step 1:} write the contrapositive: if $m$ is even or $n$ is even, then
$mn$ is even. \textbf{Step 2:} use cases. If $m=2k$, then $mn=2(kn)$; if
$n=2\ell$, then $mn=2(m\ell)$. \textbf{Step 3:} in either case the quantity in
parentheses is an integer, so $mn$ is even. The contrapositive, hence the
original claim, is true.
\end{workedbox}

\begin{workedbox}{C: no largest integer}
\textbf{Assume the opposite:} let $N$ be a largest integer. \textbf{Construct:}
$N+1$ is an integer and $N+1>N$. \textbf{Contradiction:} maximality requires
$N+1\le N$, while arithmetic gives $N+1>N$. Therefore no largest integer
exists.
\end{workedbox}

\begin{workedbox}{C: rational plus irrational}
Let $r\in\mathbb Q$ and $x\notin\mathbb Q$. \textbf{Assume the opposite:}
$r+x\in\mathbb Q$. \textbf{Isolate the irrational number:} rationals are closed
under subtraction, so $x=(r+x)-r\in\mathbb Q$. This contradicts the hypothesis
$x\notin\mathbb Q$; hence $r+x$ is irrational.
\end{workedbox}

\columnbreak

\begin{definitionbox}{Proof by cases and existence}
A \term{proof by cases} partitions all possibilities and proves the claim in
each exhaustive case. A \term{constructive existence proof} gives an explicit
witness and verifies every required property. A \term{nonconstructive proof}
establishes existence without identifying a witness.
\end{definitionbox}

\begin{workedbox}{D: cases on parity}
\textbf{Exhaustive split:} every integer $n$ has form $2k$ or $2k+1$.
\textbf{Even case:} $n(n+1)=2[k(n+1)]$. \textbf{Odd case:}
\[
n(n+1)=(2k+1)(2k+2)=2[(2k+1)(k+1)].
\]
Both bracketed quantities are integers. Since all integers occur in one case,
$n(n+1)$ is always even.
\end{workedbox}

\begin{workedbox}{D: constructive existence}
\textbf{Witnesses:} $24,25,26,27$ are consecutive and positive.
\textbf{Verification:}
\[
24=2\cdot12,\quad25=5\cdot5,\quad26=2\cdot13,\quad27=3\cdot9.
\]
Each factorization is nontrivial, so all four are composite. Listing numbers
without checking the required property would not finish the proof.
\end{workedbox}

\begin{workedbox}{D: counterexamples}
\textbf{Irrational sum claim:} choose $\sqrt2$ and $-\sqrt2$ (both irrational).
Their sum is $0$ (rational), so the hypothesis holds and the conclusion fails.

\textbf{Prime-output claim:} choose the allowed input $n=40$:
$40^2+40+41=41^2$, which is composite. One checked counterexample refutes a
universal claim.
\end{workedbox}

\begin{warningbox}{Common errors}
Cases must be exhaustive. A counterexample must obey every domain restriction.
In contradiction, do not merely restate the assumed negation; derive a genuine
conflict. In contraposition, negate both sides correctly and reverse the arrow.
\end{warningbox}

\begin{checklistbox}{C--D readiness}
Can you state a contrapositive, run a contradiction, split an exhaustive case,
verify a constructive witness, and refute a universal claim with a complete
counterexample?
\end{checklistbox}
\end{multicols}

\clearpage
\section*{E. Strategy Choice and Proof Repair}
\begin{multicols}{2}
\begin{methodbox}{Choose a promising method}
\begin{tabularx}{\linewidth}{@{}lX@{}}
Direct & Definitions turn the hypothesis into the target form.\\
Contrapositive & The negated conclusion has a useful structural definition.\\
Contradiction & The conclusion denies existence or involves irrationality.\\
Cases & The domain naturally partitions (parity, sign, membership).\\
Constructive & The claim asks only that an example exists.\\
Counterexample & The universal claim looks too strong.
\end{tabularx}
State the first substantive step, not merely the method name.
\end{methodbox}

\begin{workedbox}{E: repair reused witnesses}
\textbf{Gap:} ``$a=2k$ and $b=2k$'' forces $a=b$; a shared witness was not
justified. \textbf{Repair:} choose $j,k\in\mathbb Z$ with $a=2j$ and $b=2k$.
Then $a+b=2(j+k)$, and $j+k\in\mathbb Z$, so the sum is even.
\end{workedbox}

\begin{workedbox}{E: repair division by zero}
\textbf{Gap:} the original claim fails for $r=0$, and dividing by $r$ may be
illegal. \textbf{Correct statement:} if $r\in\mathbb Q\setminus\{0\}$ and $x$
is irrational, then $rx$ is irrational. \textbf{Proof:} if $rx$ were rational,
division by nonzero rational $r$ would give $x=(rx)/r\in\mathbb Q$, a
contradiction.
\end{workedbox}

\begin{workedbox}{E: repair circular reasoning}
\textbf{Gap:} assuming $n^2$ is odd assumes the conclusion. \textbf{Repair:}
start with the hypothesis $n=2k+1$. Then
\[
n^2=4k^2+4k+1=2(2k^2+2k)+1,
\]
The parenthesized quantity is an integer, so the result has form
$2(\text{integer})+1$ and $n^2$ is odd.
\end{workedbox}

\columnbreak

\begin{workedbox}{E: repair a bad starting element}
\textbf{Gap:} starting with $x\in A\cap B$ does not cover arbitrary $x\in A$.
\textbf{Repair:} let arbitrary $x\in A$. The equality $A=A\cap B$ gives
$x\in A\cap B$, so $x\in B$. Since $x$ was arbitrary, $A\subseteq B$.
\end{workedbox}

\begin{definitionbox}{Essential versus minor errors}
An \term{essential error} breaks the proof: assuming the conclusion, reversing
an implication, using an unjustified common witness, dividing by a possible
zero, omitting a case, or giving an example instead of a universal proof.
Minor notation is repairable when scope and reasoning remain unambiguous.
\end{definitionbox}

\begin{methodbox}{A proof-repair checklist}
Ask:
\begin{enumerate}
\item Did the proof start with the correct arbitrary object or hypothesis?
\item Are every witness and variable scoped correctly?
\item Is any implication reversed or invalid inference used?
\item Is division known to be legal?
\item Are cases exhaustive and conclusions exact?
\item Does the final line invoke the requested definition?
\end{enumerate}
Name the gap before replacing it.
\end{methodbox}

\begin{warningbox}{Presentation matters}
A proof is a readable argument, not a stack of equations. State assumptions,
connect equations with words, identify the witness produced, and finish with
the exact conclusion. A correct idea with a missing essential justification is
not a complete proof.
\end{warningbox}

\begin{checklistbox}{Whole-Core readiness}
Can you perform an inference, write a direct proof, use contraposition or
contradiction, handle cases/existence/counterexamples, and select or repair a
strategy? A live 20-minute form selects one complete problem from each section
A--E; four fully correct problems earn Meets and five earn Exceeds.
\end{checklistbox}
\end{multicols}

\end{document}
